% Imporditud: tools/import_eksperimendid.py
% Allikas: Eesti füüsikaolümpiaadi eksperimendi ülesannete
%   kogu (Rael Kalda, 2023), ülesanne Ü174, lahendus L174.
\setAuthor{Ardi Loot}
\setRound{lõppvoor}
\setYear{2017}
\setNumber{G E2}
\setDifficulty{5}
\setTopic{gaasid}
\setVahendid{õhupall, millimeetripaber, marker, teadaoleva massiga koormis}
\setPunktid{12}
\setAllikas{Eksperimendi kogu Ü174, lahendus lk 133}
\setLahendusseis{toores}

\prob{Õhupall 2}
Määrata võimalikult täpselt mitme protsendi võrra on õhupallis olev rõhk suurem

õhurõhust $p_0$ = 101 kPa. Võib lähtuda eeldusest, et palli väikeste deformatsioonide korral muutub palli kuju, mitte aga tema ruumala. Raskuskiirendus g = 9,8 m/s$^2$.

\hint

\solu
Asetame õhupalli millimeeterpaberile ja sellele koormise. Võttes eelduseks, et väikeste deformatsioonide korral ei muutu oluliselt õhupalli ruumala, siis järelikult ei muutu oluliselt ka õhupallis olev rõhk. Koormise tõttu suureneb õhupalli ja millimeetripaberi kokkupuutepind, mille leiame märkides selle markeriga paberile. Teine variant on õhupall eelnevalt markeriga kokku määrida ja mõõta paberile jääva jälje pindala. Õhupalli horisontaalseks surutud pind mõjub lauale jõuga $\Delta$pS, kus $\Delta$p on õhupallis olev lisarõhk ja S on õhupalli ja laua kokkupuutepinna suurus. Lugedes õhupalli massi tühiseks, peab see jõud olema võrdne koormise raskusjõuga mg, kus m on koormise mass. Õhupallis olev lisarõhk avaldub seega $\Delta$p = mg/S ja kokkuvõttes on õhupalli sees olev rõhk õhurõhust $\eta$ = ($p_0$ + $\Delta$p) /$p_0$ korda suurem.

Märkus: Õhupallis oleva pinge saame jätta arvestamata, kunas õhupalli ja millimeetripaberi kokkupuutepind on horisontaalne ning seega puudub pinge vertikaalne komponent.
\probend
